\clearpage
\section{The actual enlarged rank-four inducing carrier}\label{app:rank4}


\hypertarget{rank4-result-and-boundary}{%
\subsection{Result and boundary}\label{rank4-result-and-boundary}}

For the actual rank-four Artin--Schreier candidate of Block167, one can
enlarge the compact group J1 to an actual compact subgroup P with

\begin{Verbatim}[fontsize=\small,breaklines=true,breakanywhere=true]
J1 normal in P,  [P:J1]=2,
E=Ind_(J1)^P(eta) irreducible of dimension 4.
\end{Verbatim}

The representation E occurs inside the restriction of the same actual
supercuspidal candidate pi. There is an explicit normal subgroup H0 of P
with {[}P:H0{]}=16 such that P/H0 is elementary abelian and the
commutator form supplied by the actual character on H0 is nondegenerate.
Thus the carrier has an actual four-dimensional binary symplectic
quotient. No abstract transfer of matrix operators is used to create
this action.

The finite image of P has scalar center theta(H0), which can be larger
than \{+1,-1\}. It is literally an extraspecial group if this scalar
center has order two. That additional assertion is not assumed or proved
here. The construction supplies an actual finite Heisenberg central
extension with nondegenerate projective quotient F\_2\^{}4. Matching an
independently prescribed quadratic phase is also a separate assertion.

\hypertarget{rank4-actual-arithmetic-groups-and-notation}{%
\subsection{1. Actual arithmetic groups and
notation}\label{rank4-actual-arithmetic-groups-and-notation}}

Let F=F\_2((t)), and take the audited rank-four constructor with odd
0\textless h\_1\textless h\_2. Write

\begin{Verbatim}[fontsize=\small,breaklines=true,breakanywhere=true]
q_i=3h_i,  n=2q_2,  a=n/2=q_2,
r=q_1=2m+1,  m=(r-1)/2>=1.
\end{Verbatim}

We have a\textgreater r, so n\textgreater2r. Let A be the complete
period-four hereditary order, with radical P\_A. Let gamma be the actual
degree-two outer minimal approximation, E=F{[}gamma{]},
B=End\_E(F\^{}4)=M\_2(E), B0=A cap B, and Q=rad(B0). The actual final
parameter is beta=gamma+delta, and

\begin{Verbatim}[fontsize=\small,breaklines=true,breakanywhere=true]
s(delta)=c,
\end{Verbatim}

where c is the minimal degree-two Artin--Schreier parameter over E, of
B0-level r. The final field K=F{[}beta{]} is totally ramified of degree
four; no embedding E subset K is asserted.

Use \(U_A^j=1+P_A^j\) and \(U_B^j=1+Q^j\).
BK3.1.7--9 give the actual groups

\begin{Verbatim}[fontsize=\small,breaklines=true,breakanywhere=true]
H=H1(beta)=U_K^1 U_B^(m+1) U_A^(a+1),
J=J1(beta)=U_K^1 U_B^(m+1) U_A^a.
\end{Verbatim}

In particular J/H has dimension two over F\_2. Let theta be the actual
full simple character constructed for beta, and let eta be its unique
Heisenberg representation on J, of dimension two. The argument below
works with this supplied actual theta; it does not replace its values.

Choose the outer minimal simple character theta\_gamma from its defining
sequence. On H cap U\_A\^{}(m+1), the recursion says

\begin{Verbatim}[fontsize=\small,breaklines=true,breakanywhere=true]
theta=theta_gamma psi_delta.
\end{Verbatim}

Moreover H is contained in H1(gamma)=U\_B\^{}1 U\_A\^{}(a+1). The
character theta\_gamma is invariant under U\_B\^{}1: on U\_B\^{}1 it
factors through the E-determinant, and on U\_A\^{}(a+1) it is the
gamma-additive character.

\hypertarget{rank4-the-actual-final-field-has-the-needed-graded-line}{%
\subsection{2. The actual final field has the needed graded
line}\label{rank4-the-actual-final-field-has-the-needed-graded-line}}

For every integer j\textgreater=1 and k in K cap P\_A\^{}j, the
minimal-approximation commutator estimate gives

\begin{Verbatim}[fontsize=\small,breaklines=true,breakanywhere=true]
k=b+e,  b in Q^j,  e in P_A^(j+n-r).              (1)
\end{Verbatim}

Indeed {[}gamma,k{]}=-{[}delta,k{]} belongs to P\_A\^{}(j-r), and the
minimal corestriction exact sequence for gamma solves its commutator
with an error of grade j+n-r. The remaining difference centralizes gamma
and therefore belongs to B. This is the usual minimal case of BK1.4.10.

Applying s to {[}beta,k{]}=0 and using (1) gives

\begin{Verbatim}[fontsize=\small,breaklines=true,breakanywhere=true]
[c,b] in Q^(j+n-2r).
\end{Verbatim}

Hence the leading B-grade of b commutes with the leading grade of c. Let
R be the complete period-two E-chain shift, with R\^{}2=z I, where z is
the chosen uniformizer of E. In the constructor's basis

\begin{Verbatim}[fontsize=\small,breaklines=true,breakanywhere=true]
c congruent to R^(-r) modulo Q^(1-r).
\end{Verbatim}

Because r is odd, the centralizing subspace on every B-grade j is the
one-dimensional line spanned by R\^{}j. On the other hand K is a totally
ramified maximal field normalizing the complete A-chain, so (K cap
P\_A\^{}j)/(K cap P\_A\^{}(j+1)) has dimension one over F\_2. Estimate
(1) is injective on the leading grade, so its image is exactly that
centralizing line.

This proves the required graded assertion for the actual K, without
replacing it by the derived model E{[}c{]} or assuming nested fields.

\hypertarget{rank4-normality-and-index-of-the-enlarged-group}{%
\subsection{3. Normality and index of the enlarged
group}\label{rank4-normality-and-index-of-the-enlarged-group}}

Set

\begin{Verbatim}[fontsize=\small,breaklines=true,breakanywhere=true]
P=J U_B^m.
\end{Verbatim}

We first prove this is a group. The group U\_B\^{}m normalizes
U\_B\^{}(m+1) and every U\_A\^{}j. For k in U\_K\^{}1, (1) approximates
k by an element of U\_B\^{}1 with error in U\_A\^{}(n-r+1). Therefore

\begin{Verbatim}[fontsize=\small,breaklines=true,breakanywhere=true]
[U_B^m,U_K^1]
   subset U_B^(m+1) U_A^(m+n-r+1)
   subset H,
\end{Verbatim}

where the last inclusion follows from a\textgreater r=2m+1. Thus
U\_B\^{}m normalizes both H and J, and P is an actual compact subgroup
with J normal in P.

The quotient \(U_B^m/U_B^{m+1}\) has dimension two. By
Section~2 the image of J cap U\_B\^{}m on this grade is exactly its
centralizing line, of dimension one. For the converse inclusion needed
here, the B-approximation of a K-element of grade m differs from it by
grade m+n-r\textgreater a; that error lies in U\_A\^{}a subset J. Thus
the whole centralizing line really occurs in the intersection, not only
in an abstract associated graded space. Consequently

\begin{Verbatim}[fontsize=\small,breaklines=true,breakanywhere=true]
[P:J]=[U_B^m:J cap U_B^m]=2.                    (2)
\end{Verbatim}

For a concrete representative choose D=diag(1,0) in B and put

\begin{Verbatim}[fontsize=\small,breaklines=true,breakanywhere=true]
x=R^m D,  u=1+x.
\end{Verbatim}

Its leading grade is outside the centralizing line, so u represents the
nontrivial coset of P/J. Since \(u^2=1+x^2\) in
characteristic two, we have u\^{}2 in U\_B\^{}(2m) subset U\_B\^{}(m+1)
subset H.

\hypertarget{rank4-the-actual-character-has-an-orbit-of-length-two}{%
\subsection{4. The actual character has an orbit of length
two}\label{rank4-the-actual-character-has-an-orbit-of-length-two}}

Define the character of H

\begin{Verbatim}[fontsize=\small,breaklines=true,breakanywhere=true]
lambda(h)=theta(u h u^(-1))/theta(h).
\end{Verbatim}

For h in H the preceding normality argument gives {[}u,h{]} in H cap
U\_A\^{}(m+1). Since theta\_gamma is invariant under u,

\begin{Verbatim}[fontsize=\small,breaklines=true,breakanywhere=true]
lambda(h)=psi_delta([u,h]) in {+1,-1}.          (3)
\end{Verbatim}

Thus lambda has order dividing two, and its values are actual values of
the supplied simple character. It is J-invariant, since both theta and
its u-conjugate are J-invariant.

It is nontrivial. Take y=R\^{}(m+1)D and h=1+y in U\_B\^{}(m+1). The
leading commutator in (3), after corestriction, has value

\begin{Verbatim}[fontsize=\small,breaklines=true,breakanywhere=true]
psi_E(tr_(M_2(E)/E)(c[x,y]))=-1.
\end{Verbatim}

Here the commutator remainder has B-grade at least 3m+1, which becomes
positive after multiplication by c of grade -2m-1. On the leading
grades, writing X=R\^{}m diag(A,B) and Y=R\^{}(m+1)diag(C,D), one
obtains

\begin{Verbatim}[fontsize=\small,breaklines=true,breakanywhere=true]
tr(R^(-2m-1)[X,Y])=(A+B)(C+D) in F_2.
\end{Verbatim}

Both sums equal one for the chosen D=diag(1,0). The normalized psi\_E is
nontrivial on the residue field and trivial on its maximal ideal. For
this AS normalization one can check the sign directly from
psi\_E(1)=psi\_F(Tr\_(E/F)(alpha))=psi\_F(1)=-1.

Hence theta\^{}u differs from theta, eta\^{}u differs from eta, and
Clifford theory applied to (2) gives

\begin{Verbatim}[fontsize=\small,breaklines=true,breakanywhere=true]
E_carrier=Ind_J^P(eta) irreducible, dim(E_carrier)=4.       (4)
\end{Verbatim}

\hypertarget{rank4-an-actual-nondegenerate-symplectic-quotient-of-dimension-four}{%
\subsection{5. An actual nondegenerate symplectic quotient of dimension
four}\label{rank4-an-actual-nondegenerate-symplectic-quotient-of-dimension-four}}

Put H0=ker(lambda) inside H, and theta0=theta restricted to H0. In the
two-coset induction model, h in H acts by

\begin{Verbatim}[fontsize=\small,breaklines=true,breakanywhere=true]
diag(theta(h) I_2, theta^u(h) I_2).
\end{Verbatim}

Thus it acts scalarly precisely when h belongs to H0. We check that H0
is the entire scalar preimage and that its quotient is symplectic.

First, lambda is invariant under u. Indeed u\^{}2 belongs to J and theta
is J-invariant, so \(\lambda^u=\lambda^{-1}=\lambda\) by (3).
Consequently H0 is normal in P, and theta0 is P-invariant.

Second, lambda is trivial on U\_A\^{}(a+1). There the correction formula
is valid, gamma commutes with u, and the delta contribution to a
commutator has grade at least a+1+m-r=a-m\textgreater0. Therefore it
lies in the kernel of the additive character.

Write J=H U\_A\^{}a. Since lambda is a character and J-invariant, it
kills {[}H,H{]} and {[}H,U\_A\^{}a{]}. It also kills
\([U_A^a,U_A^a]\), which lies in U\_A\^{}(2a). Hence
lambda({[}J,J{]})=1. It kills every square in J as well: squares of H
have lambda-value one, squares of U\_A\^{}a lie in U\_A\^{}(2a), and the
mixed commutator contributions have just been handled.

Finally lambda({[}u,H{]})=1 by u-invariance, and {[}u,U\_A\^{}a{]}
belongs to U\_A\^{}(a+m) subset U\_A\^{}(a+1), so lambda({[}u,J{]})=1.
Also lambda(u\^{}2)=1 because u commutes with u\^{}2. It follows that

\begin{Verbatim}[fontsize=\small,breaklines=true,breakanywhere=true]
[P,P] subset H0,  g^2 in H0 for every g in P,
[P:H0]=[P:J][J:H][H:H0]=2*4*2=16.
\end{Verbatim}

Therefore P/H0 is an actual elementary abelian group of dimension four.
If an element of J acts scalarly on the carrier, its first eta-block is
scalar, so it belongs to H by the original nondegenerate Heisenberg
quotient J/H. It then belongs to H0. An element outside J exchanges the
two coset summands and cannot act scalarly. Thus H0 is precisely the
scalar preimage.

The pairing

\begin{Verbatim}[fontsize=\small,breaklines=true,breakanywhere=true]
(gH0,kH0) |-> theta0([g,k])
\end{Verbatim}

is alternating and well defined. Since the projective quotient has
exponent two, its values lie in \{+1,-1\}. It is nondegenerate: an
element in its radical commutes with the entire irreducible carrier
image and therefore acts scalarly by Schur's lemma, forcing it into H0.

This is the actual four-dimensional binary Heisenberg quotient sought
for capacity. Its operators are supplied by elements of P.

\hypertarget{rank4-the-carrier-occurs-in-the-same-supercuspidal-candidate}{%
\subsection{6. The carrier occurs in the same supercuspidal
candidate}\label{rank4-the-carrier-occurs-in-the-same-supercuspidal-candidate}}

Let pi=c-Ind\_N\^{}(GL\_4(F))(Lambda) be the existing actual candidate,
with N=K\^{}times J. The residue field is F\_2, so the compact extended
type restricts to eta on J. Since P is contained in U\_A\^{}1,

\begin{Verbatim}[fontsize=\small,breaklines=true,breakanywhere=true]
N cap P=J.
\end{Verbatim}

Indeed an element k j in P with k in K\^{}times has k in U\_A\^{}1,
hence k in U\_K\^{}1 subset J. The identity double-coset term in the
restriction of compact induction to P is therefore

\begin{Verbatim}[fontsize=\small,breaklines=true,breakanywhere=true]
Ind_J^P(eta)=E_carrier.
\end{Verbatim}

It is realized concretely by the functions supported on N P. Thus the
four-dimensional carrier occurs in the same pi, with its already fixed
full type character and noncompact extension scalar. No new arithmetic
candidate, beta extension, or inverse correspondence is used.

\hypertarget{rank4-exact-limits-finite-image-versus-literal-extraspecial-group}{%
\subsection{7. Exact limits: finite image versus literal extraspecial
group}\label{rank4-exact-limits-finite-image-versus-literal-extraspecial-group}}

The actual finite image is P/ker(theta0), because H0 is the scalar
preimage. Its center is the scalar cyclic 2-group theta0(H0), and its
commutator subgroup has order two. Its quotient by its center is
F\_2\^{}4.

If theta0(H0) has order two, this finite image is extraspecial of order
32. If its scalar center is larger, quotienting arbitrary central
elements need not preserve its commutator: in a cyclic 2-group the
order-two subgroup can lie inside the squares. Thus a literal
extraspecial quotient is not justified solely by the calculations above.

The stronger statement proved is the one with full actual character: an
actual compact subgroup, actual induced carrier, actual scalar subgroup
and nondegenerate actual commutator pairing. Passing to the usual full
scalar pushout retains all theta0-values. It does not prove that an
independently supplied quadratic refinement, stabilizer action, or
source phase agrees with this carrier. Those requirements must be
checked separately. No H5 or universal claim~66 closure is asserted.

\hypertarget{rank4-provenance}{%
\subsection{Provenance}\label{rank4-provenance}}

Uses the audited dyadic parameter/character construction, BK1.4.10,
BK3.1.7--9, and the defining-sequence character formula. The source
normalization and all field/chain data are the actual ones in
\path{dyadic167/tower/DYADIC_AS_TOWER167.md}. This note changes neither that
proof nor its recorded exact computation. Clifford induction and the
identity compact-induction double-coset term are used in their ordinary
finite-index forms, with hypotheses proved above.
